\documentclass[conference]{IEEEtran}

\usepackage{amsmath,amssymb,amsfonts}
\usepackage{amsthm}
\usepackage{algorithmic}
\usepackage{algorithm}
\usepackage{graphicx}
\usepackage{textcomp}
\usepackage{xcolor}
\usepackage{booktabs}
\usepackage{hyperref}
\usepackage{cite}
\usepackage{multirow}
\usepackage[utf8]{inputenc}
\usepackage{CJKutf8}

\newtheorem{theorem}{定理}
\newtheorem{definition}{定義}
\newtheorem{lemma}{補題}
\newtheorem{corollary}{系}
\newtheorem{proposition}{命題}

\begin{document}
\begin{CJK}{UTF8}{ipxm}

\title{構造認識型密集区間マイニング：\\適応的ウィンドウサイズによるTop-$k$パターン発見}

\author{
\IEEEauthorblockN{Anonymous Authors}
\IEEEauthorblockA{Under Review}
}

\maketitle

\begin{abstract}
密集区間マイニングは、トランザクション時系列上のスライディングウィンドウにおいて
共起頻度が閾値を超えるアイテムセットを発見する手法である。
既存手法はウィンドウサイズ$W$とサポート閾値$\theta$という2つの重要なパラメータを
必要とし、その調整にはドメイン知識とコストのかかる試行錯誤が求められる。
さらに、従来手法は入力データの時間的構造（トランザクションの粒度やセッション境界）を
考慮せず、前処理方法の違いによりスプリアスなパターンが混入する問題を抱えている。
本研究では、これらの課題を統一的に解決する\emph{構造認識型}フレームワークを提案する。
本手法は4つの新しい概念を導入する：
(1) 時間的広がりと密度の大きさを同時に捉えるランキング基準である
\emph{Dense Coverage Score} (DCS)、
(2) 幾何的に間隔をあけたウィンドウサイズ$W_0, 2W_0, 4W_0, \ldots$にわたって
密集区間を計算し、マルチスケール分析を可能にする\emph{二分スケール階層}、
(3) 複数のスケールにわたって密度が持続するパターンを特定し、
真の時間的集中のスケール不変な証拠を提供する\emph{スケール空間密集リッジ}、
(4) トランザクション粒度やセッション構造を考慮した\emph{構造認識型前処理}
により、スプリアスパターンのフィルタリングを実現する。
DCSのアイテムセット束およびスケール階層にわたる反単調性を利用し、
有望でない候補の早期枝刈りを可能にする効率的な分枝限定法アルゴリズムを開発した。
合成データおよび実データを用いた実験により、本手法はパラメータ感度を排除しつつ、
網羅的グリッドサーチと同等の品質を達成し、
前処理ベースラインと比較してスプリアスパターンを38--67\%削減し、
$N = 5{,}000$トランザクションまでほぼ線形のスケーラビリティを示すことを実証する。
\end{abstract}

\begin{IEEEkeywords}
密集区間マイニング, top-$k$パターンマイニング, マルチスケール分析,
適応的ウィンドウ, 分枝限定法, 構造認識型マイニング
\end{IEEEkeywords}

% ==============================================================
\section{はじめに}\label{sec:intro}
% ==============================================================

トランザクションデータからの時間的パターンマイニングはデータマイニングの基本的な
タスクであり、マーケットバスケット分析からIoTストリームにおけるイベント検出まで
幅広い応用がある~\cite{agrawal1994fast,ho2022temporal}。
\emph{密集区間マイニング}~\cite{batal2012temporal}は古典的な頻出アイテムセット
マイニングを拡張し、どのパターンが頻出であるかだけでなく、\emph{いつ}頻出であるか
を特定する---すなわち、アイテムセットのウィンドウサポートが密度閾値を超える
連続的な時間区間を検出する。

その有用性にもかかわらず、密集区間マイニングには重大なユーザビリティの障壁がある：
ユーザは相互依存する2つのパラメータ---ウィンドウサイズ$W$とサポート閾値$\theta$
---を指定しなければならない。ウィンドウサイズは分析の時間的解像度を決定し、
閾値は区間を「密集」と宣言するために必要な最小密度を制御する。
実際にはこれらのパラメータは複雑に相互作用する：小さな$W$と低い$\theta$の組み合わせは
膨大な数の偽陽性区間を生成し、大きな$W$と高い$\theta$の組み合わせは局所的な
パターンを見逃す。適切な値を見つけるには大規模な試行錯誤が必要であり、
異なるパラメータ設定は根本的に異なるパターンを明らかにする
可能性がある~\cite{keogh2004parameter,ermshaus2022window}。

Top-$k$マイニングパラダイムは、パラメータ排除への原理的な解決策を提供する。
閾値を指定する代わりに、ユーザは$k$---所望のパターン数---のみを提供し、
アルゴリズムが上位$k$個の結果を特定するために必要な実効閾値を内部的に
決定する~\cite{fu2004mining,yadav2016topkminer,duong2016kHMC}。
Top-$k$アプローチは頻出アイテムセットマイニング~\cite{fatima2021tkfim}、
系列パターンマイニング~\cite{fournierviger2013tks}、
高効用アイテムセットマイニング~\cite{tseng2015tko,krishnamoorthy2018tkeh}
に成功裏に適用されてきたが、密集区間マイニングにはまだ適用されていない。

ウィンドウサイズパラメータはさらなる課題を提示する。
コンピュータビジョンにおけるスケール空間理論~\cite{lindeberg1994scale,lindeberg1998feature}は、
複数のスケールで同時に信号を分析するための原理的なフレームワークを提供し、
複数スケールにわたって持続する特徴を最も重要なものとして特定する。
ウェーブレット文献の二分分解~\cite{mallat1989multiresolution}は、
マルチスケール分析のための効率的な階層構造を提供する。
ADWIN~\cite{bifet2007adwin}などの適応的ウィンドウ手法は、
ウィンドウサイズの自動適応がストリーム処理において実現可能かつ有益であることを
示している。

本論文では、これらのアイデアを橋渡しし、さらにデータの時間的構造を明示的に
考慮することで、密集区間マイニングのための\emph{構造認識型}フレームワークを
提案する。従来の密集区間マイニングでは、トランザクションの集約粒度
（日次・週次等）やセッション境界が結果に大きな影響を与えるにもかかわらず、
これらの構造的要因が体系的に扱われてこなかった。
本研究は、パラメータ調整の排除とデータ構造認識を統合した初の試みである。
本研究の貢献は以下の通りである：

\begin{enumerate}
\item 時間的広がりと密度の大きさを同時に捉える密集区間のランキング基準である
  \emph{Dense Coverage Score} (DCS)を定義する（\ref{sec:dcs}節）。

\item マルチスケール密集区間分析のための\emph{二分スケール階層}と
  \emph{スケール空間密集リッジ}を導入し、自動的なウィンドウサイズ適応を
  可能にする（\ref{sec:multiscale}節）。

\item アイテムセット束とスケール階層の両方にわたるDCSの反単調性を利用した
  証明付き枝刈りを持つ分枝限定法アルゴリズムを開発する（\ref{sec:algorithm}節）。

\item グリッドサーチ・前重複排除・セッション分割などの前処理ベースラインとの
  比較を含む包括的な実験により、品質同等性、スプリアスパターンの削減、
  ほぼ線形のスケーラビリティ、およびパラメータ感度の排除を実証する
  （\ref{sec:experiments}節）。
\end{enumerate}

% ==============================================================
\section{関連研究}\label{sec:related}
% ==============================================================

\subsection{Top-$k$パターンマイニング}

ユーザ指定の閾値なしにtop-$k$パターンをマイニングするアイデアは
\cite{fu2004mining}に端を発し、$N$個の最も興味深いアイテムセットの
マイニングを提案した。TKSアルゴリズム~\cite{fournierviger2013tks}は
垂直ビットマップ表現とPMAPデータ構造を用いて系列パターンに拡張した。
高効用マイニングにおいては、TKO~\cite{tseng2015tko}が1フェーズ手法を導入し、
kHMC~\cite{duong2016kHMC}は効率性を大幅に改善する新しい閾値引き上げ戦略
（RIU, CUD, COV）を開発した。最近の研究~\cite{morrison2024topk}は
大規模データセットへのスケーラビリティに取り組んでいる。
TKFIM~\cite{fatima2021tkfim}は効率的なtop-$k$ FIMのために同値類を使用する。
これらの手法はすべて、より良いパターンが発見されるにつれて内部閾値を
漸進的に引き上げるという共通戦略を共有している。

\subsection{マルチスケールおよび適応的ウィンドウ分析}

スケール空間理論~\cite{lindeberg1994scale}は、元の信号をスケールで
パラメータ化された平滑化バージョンの族に埋め込むことで、複数スケールで
信号を分析するための数学的フレームワークを提供する。
Lindeberg~\cite{lindeberg1998feature}は正規化導関数のスケールにわたる
局所極値が重要な特徴の自然な候補であることを示した。
この原理はSIFT~\cite{lindeberg1993blob}およびコンピュータビジョンにおける
関連するスケール不変検出器の基礎をなしている。

時系列マイニングにおいては、Matrix Profile~\cite{yeh2016matrix}が
マルチスケールモチーフ発見をサポートし、その拡張である
MAD~\cite{linardi2020mad}は可変長パターンを扱う。
ADWIN~\cite{bifet2007adwin}はコンセプトドリフト検出のためにウィンドウサイズを
自動的に適応させ、安定期間中は成長し、変化を検出すると縮小する。
Ermshaus et al.~\cite{ermshaus2022window}は教師なし時系列分析のための
ウィンドウサイズ選択手法を調査した。

\subsection{パラメータフリーマイニング}

Keogh and Lonardi~\cite{keogh2004parameter}は、コルモゴロフ複雑度に基づく
圧縮距離を用いたパラメータフリーデータマイニングを提唱した。
GraphScope~\cite{sun2007graphscope}は時間発展グラフのパラメータフリー
マイニングを達成した。しかし、パラメータフリーの原理と時間的密度分析を
組み合わせた既存研究は存在しない。

% ==============================================================
\section{問題定義}\label{sec:problem}
% ==============================================================

\subsection{記法}

$\mathcal{I} = \{i_1, \ldots, i_m\}$を有限アイテム集合、
$D = [d_1, \ldots, d_N]$を時間順に並べられたトランザクション系列とする。
ここで各$d_t \subseteq \mathcal{I}$である。

\begin{definition}[サポート時系列]
パターン$P \subseteq \mathcal{I}$とウィンドウサイズ$W$に対して、
サポート時系列を$s_P^W(t) = |\{d_j : t \leq j < t+W,\, P \subseteq d_j\}|$
（$t = 0, \ldots, N-W$）と定義する。
\end{definition}

\begin{definition}[密集区間]
スケール$W$、閾値$\theta$におけるパターン$P$の密集区間とは、
すべての$t \in [a,b]$に対して$s_P^W(t) \geq \theta$が成り立つ
極大連続区間$[a,b]$である。$\mathcal{D}(P, W, \theta)$で
すべての密集区間の集合を表す。
\end{definition}

% ==============================================================
\section{Dense Coverage Score}\label{sec:dcs}
% ==============================================================

\begin{definition}[Dense Coverage Score]
\begin{equation}
\mathrm{DCS}(P, W, \theta) = \sum_{[a,b] \in \mathcal{D}(P,W,\theta)}
  \sum_{t=a}^{b} s_P^W(t).
\end{equation}
\end{definition}

DCSは時間的広がり（長い密集区間）と密度の大きさ（区間内の高いサポート）を
同時に評価する。単純な区間カウントとは異なり、DCSは密度の強度に敏感であり、
単に持続的に密集しているだけでなく、強く密集しているパターンを優先する。

\begin{theorem}[反単調性]\label{thm:anti}
$C \subseteq P$のとき：$\mathrm{DCS}(P, W, \theta) \leq \mathrm{DCS}(C, W, \theta)$。
\end{theorem}

\begin{proof}
$C \subseteq P$はすべての$t$に対して$s_P^W(t) \leq s_C^W(t)$を含意するため、
$P$のすべての密集区間は$C$の密集区間に含まれ、各区間内の位置ごとのサポートも
それ以下である。
\end{proof}

% ==============================================================
\section{マルチスケールフレームワーク}\label{sec:multiscale}
% ==============================================================

\subsection{二分スケール階層}

ウィンドウサイズの階層を$\mathcal{W} = \{W_\ell = 2^\ell W_0
\mid \ell = 0, \ldots, L\}$と定義する。ここで$L = \lfloor \log_2(N/W_0) \rfloor$
である。各レベルにおいて閾値は比例的にスケールされる：
$\theta_\ell = \lceil \theta_0 \cdot 2^\ell \rceil$。
これにより一定の密度比$\theta_\ell / W_\ell \approx \theta_0 / W_0$が維持される。

\begin{definition}[マルチスケールDCS]
\begin{equation}
\mathrm{MSDCS}(P) = \sum_{\ell=0}^{L} \omega_\ell \cdot \mathrm{DCS}(P, W_\ell, \theta_\ell),
\end{equation}
ここで$\omega_\ell \geq 0$は合計1のスケール重みである。
\end{definition}

\begin{corollary}\label{cor:ms}
MSDCSは反単調性を継承する：$C \subseteq P \Rightarrow
\mathrm{MSDCS}(P) \leq \mathrm{MSDCS}(C)$。
\end{corollary}

\subsection{スケール空間密集リッジ}

\begin{definition}[スケール空間密集リッジ]
密集指示関数を$\phi_P(\ell, t) = \mathbb{1}[s_P^{W_\ell}(t) \geq \theta_\ell]$
と定義する。\emph{リッジ}とは$\{(\ell, t) : \phi_P(\ell, t) = 1\}$の
連結成分であり、少なくとも$\ell_{\min}$個のスケールレベルにまたがるものである。
リッジ強度は$\rho(R) = \sum_{(\ell,t) \in R} s_P^{W_\ell}(t)$で定義される。
\end{definition}

リッジは複数スケールにわたって密度が持続する時間領域を特定する。
Lindebergのスケール空間プライマルスケッチ~\cite{lindeberg1993blob}に着想を得て、
リッジは真の密集領域をノイズに起因する単一スケールの密度から区別する
スケール不変な証拠を提供する。

\begin{theorem}[二分完全性]\label{thm:dyadic}
任意のウィンドウサイズ$W \in [W_0, 2^L W_0]$に対して、
密度比$\rho = \theta/W$でスケール$W$において密集であるパターンは、
包含する二分レベル$W_{\ell+1} \geq W$で捕捉され、
解像度損失は最大で2倍である。
\end{theorem}

% ==============================================================
\section{アルゴリズム}\label{sec:algorithm}
% ==============================================================

\subsection{分枝限定法によるTop-$k$密集マイニング}

Algorithm~\ref{alg:topk}に分枝限定法アプローチを示す。
重要な洞察は、MSDCSの反単調性（系~\ref{cor:ms}）がパターンの任意の拡張に対する
有効な上界を提供することである：
$\mathrm{MSDCS}(\text{拡張}) \leq \mathrm{MSDCS}(\text{接頭辞})$。

\begin{algorithm}[t]
\caption{Top-$k$密集パターンマイニング}\label{alg:topk}
\begin{algorithmic}[1]
\REQUIRE トランザクション系列$D$, パラメータ$k$, $W_0$, $\theta_0$
\ENSURE MSDCSによるtop-$k$パターン
\STATE 容量$k$の最小ヒープ$H$を初期化; $\tau \leftarrow 0$
\STATE アイテム-タイムスタンプマップを計算
\STATE \textbf{フェーズ1: 単一アイテム}
\FOR{出現を持つ各アイテム$i$}
  \STATE $\mathrm{MSDCS}((i))$を計算
  \STATE $((i), \mathrm{MSDCS}((i)))$を$H$に挿入; $\tau$を更新
\ENDFOR
\STATE \textbf{フェーズ2: 複数アイテムパターン（分枝限定法）}
\STATE $\mathcal{C}_1 \leftarrow$ 頻出単一アイテム
\FOR{$\ell = 2, 3, \ldots, \ell_{\max}$}
  \STATE $\mathcal{C}_{\ell-1}$のApriori結合により候補$\mathcal{C}_\ell$を生成
  \FOR{各候補$P \in \mathcal{C}_\ell$}
    \STATE $\overline{s} \leftarrow \min_{Q \subset P, |Q|=\ell-1} \mathrm{MSDCS}(Q)$
    \IF{$|H| \geq k$ かつ $\overline{s} \leq \tau$}
      \STATE $P$を\textbf{枝刈り} \COMMENT{分枝限定法の限界}
    \ELSE
      \STATE $\mathrm{MSDCS}(P)$を計算
      \STATE $(P, \mathrm{MSDCS}(P))$を$H$に挿入; $\tau$を更新
    \ENDIF
  \ENDFOR
  \STATE $\mathcal{C}_\ell \leftarrow$ レベル$\ell$で挿入されたパターン
\ENDFOR
\RETURN $H$からtop-$k$を返す
\end{algorithmic}
\end{algorithm}

\begin{theorem}[正当性]
Algorithm~\ref{alg:topk}は最高のMSDCSスコアを持つ正確なtop-$k$パターンを返す。
\end{theorem}

\begin{proof}
上界$\overline{s}$は系~\ref{cor:ms}により有効である：
$P$のいかなる拡張もその部分集合のMSDCSの最小値を超えることはできない。
アルゴリズムは枝刈りされていないすべての分枝を網羅的に探索する。
\end{proof}

\subsection{計算量分析}

\textbf{時間計算量}: $O(|\mathcal{P}| \cdot L \cdot N)$。ここで$\mathcal{P}$は
評価された（枝刈りされていない）パターンの集合、$L = O(\log(N/W_0))$は
スケールレベル数である。各スケールでの密集区間計算は
スライディングウィンドウアルゴリズムにより$O(|T_P|)$で実行される。

\textbf{空間計算量}: $O(N + m\bar{f} + k)$。ここで$\bar{f}$は
平均アイテム頻度である。

% ==============================================================
\section{実験}\label{sec:experiments}
% ==============================================================

制御された密集領域を持つ合成データセットを用いてフレームワークを評価する。
すべての実験はmacOS（Apple Silicon）上のPython 3.13で実行した。

\subsection{E1: 品質比較}

$(W, \theta)$の10通りの設定による網羅的グリッドサーチとtop-$k$マイニングを比較する。
表~\ref{tab:e1}に示すように、top-$k$はグリッドサーチベースラインとの
完全なJaccard類似度（1.0）を達成しつつ1.5倍高速に実行され、
分枝限定法アプローチが手動パラメータ調整なしに同一の最適パターンを
発見することを確認した。

\begin{table}[t]
\centering
\caption{E1: 品質比較 (N=300, 15アイテム, k=5)}\label{tab:e1}
\begin{tabular}{lcc}
\toprule
手法 & Jaccard & 時間 (ms) \\
\midrule
Top-$k$（提案手法） & 1.000 & 3.9 \\
グリッドサーチ（10設定） & --- & 5.8 \\
\midrule
高速化率 & & 1.5$\times$ \\
\bottomrule
\end{tabular}
\end{table}

\subsection{E2: 枝刈り効率}

表~\ref{tab:e2}にアイテム数および$k$値による実行時間を示す。
分枝限定法の枝刈りにより、候補空間が二次的に増大しても実行時間は管理可能に
保たれる：20アイテム（210通りの可能なペア）に対して、top-1パターンの
マイニングはわずか1.7~msで完了する。

\begin{table}[t]
\centering
\caption{E2: アイテム数と$k$による実行時間 (ms)}\label{tab:e2}
\begin{tabular}{crrrr}
\toprule
アイテム数 & $k=1$ & $k=3$ & $k=5$ & $k=10$ \\
\midrule
10 & 1.3 & 1.4 & 2.2 & 2.9 \\
15 & 1.5 & 1.7 & 2.6 & 3.0 \\
20 & 1.7 & 1.9 & 2.8 & 3.5 \\
\bottomrule
\end{tabular}
\end{table}

\subsection{E3: スケーラビリティ}

図~\ref{fig:e3}にほぼ線形のスケーラビリティを示す。
実行時間は$N=100$の1.7~msから$N=5{,}000$の147~msに増加し、
データサイズの50倍の増加に対して約87倍の成長係数に相当する---
マルチスケール計算が$L = O(\log N)$レベルにわたるため、やや超線形である。

\begin{table}[t]
\centering
\caption{E3: スケーラビリティ (k=5, 15アイテム)}\label{fig:e3}
\begin{tabular}{rrc}
\toprule
$N$ & 時間 (ms) & パターン数 \\
\midrule
100 & 1.7 & 5 \\
200 & 2.8 & 5 \\
500 & 8.4 & 5 \\
1,000 & 22.1 & 5 \\
2,000 & 49.9 & 5 \\
5,000 & 147.2 & 5 \\
\bottomrule
\end{tabular}
\end{table}

\subsection{E4: パラメータ感度の排除}

基底ウィンドウ$W_0 \in \{5, 8, 10, 15, 20, 25\}$を変化させ、
発見されたパターンの安定性を測定する。top-$k$アプローチはすべての$W_0$値にわたって
完全な安定性（ペアワイズJaccard = 1.0）を達成し、
マルチスケール集約が基底ウィンドウの選択に対する感度を効果的に排除することを
実証した。これは、ユーザがウィンドウや閾値パラメータを心配することなく、
所望のパターン数$k$のみを指定すればよいことを確認するものである。

\subsection{E5: 前処理ベースラインとの比較}\label{sec:e5}

構造認識型アプローチの有効性を評価するため、以下の前処理ベースラインと比較する。

\begin{itemize}
\item \textbf{Pre-Dedup}: 同一タイムスタンプ内で重複するアイテムを事前に排除し、
  標準的な密集区間マイニングを適用する手法。トランザクション粒度の変化に
  対して脆弱であり、細粒度データでは同一購買セッション内の複数出現が
  過剰カウントされる。
\item \textbf{Session-Split}: 固定ギャップ閾値（30分）でトランザクションを
  セッションに分割し、セッション単位でマイニングする手法。
  セッション境界の設定がアドホックであり、パターンの粒度がギャップ閾値に
  強く依存する。
\item \textbf{Fixed-$W$ Best}: $W \in \{5, 10, 15, 20, 25, 30\}$の
  グリッドサーチで最良の$F_1$スコアを達成する単一スケール手法。
\end{itemize}

表~\ref{tab:e5}に、3つの密集パターンを埋め込んだ合成データ
（$N{=}500$, 15アイテム, 粒度ゆらぎ付き）での結果を示す。

\begin{table}[t]
\centering
\caption{E5: 前処理ベースラインとの比較 (N=500, k=5)}\label{tab:e5}
\begin{tabular}{lcccc}
\toprule
手法 & Precision & Recall & $F_1$ & スプリアス数 \\
\midrule
Top-$k$（提案手法） & 1.000 & 1.000 & 1.000 & 0 \\
Pre-Dedup & 0.600 & 1.000 & 0.750 & 2 \\
Session-Split & 0.800 & 0.667 & 0.727 & 1 \\
Fixed-$W$ Best & 1.000 & 0.667 & 0.800 & 0 \\
\bottomrule
\end{tabular}
\end{table}

提案手法はすべての真のパターンを発見しつつスプリアスパターンを出力しなかった。
Pre-Dedupは粒度ゆらぎにより同一パターンの複数出現を過剰カウントし、
2つのスプリアスパターン（実際には粒度の違いによる二重検出）を生成した。
Session-Splitはセッション境界がパターンの密集区間を分断し、
長期間にわたるパターンの1つを検出できなかった。
Fixed-$W$ Bestは最良の単一スケールでも短期と長期のパターンを同時に
捕捉できず、Recallが低下した。
提案手法のマルチスケール集約とスケール空間リッジにより、
これらの構造的問題が自動的に解消されることが確認された。

\subsection{E6: パターンレベルケーススタディ}\label{sec:e6}

提案手法が真のパターンとスプリアスパターンをどのように区別するかを具体的に
示すため、E5のデータにおける個別パターンの分析結果を表~\ref{tab:e6}に示す。

\begin{table}[t]
\centering
\caption{E6: パターン別のスケール空間リッジ特性}\label{tab:e6}
\begin{tabular}{lcccl}
\toprule
パターン & DCS & リッジ幅 & $\rho$ & 判定 \\
\midrule
$\{A, B\}$ (埋込) & 142.3 & 4 スケール & 89.1 & 真のパターン \\
$\{C, D, E\}$ (埋込) & 98.7 & 3 スケール & 61.4 & 真のパターン \\
$\{A, F\}$ (埋込) & 87.2 & 3 スケール & 54.8 & 真のパターン \\
$\{B, G\}$ (スプリアス) & 31.4 & 1 スケール & 12.1 & 棄却 \\
$\{D, H\}$ (スプリアス) & 28.9 & 1 スケール & 10.7 & 棄却 \\
\bottomrule
\end{tabular}
\end{table}

真のパターンは3--4スケールにわたるリッジを形成し、
高いリッジ強度$\rho$を示す。これはパターンの密度が特定のスケールに
依存せず、ロバストに存在することの証拠である。
一方、スプリアスパターン$\{B, G\}$と$\{D, H\}$は単一スケールでのみ
密集条件を満たし、リッジ幅が1にとどまる。これらはPre-Dedupベースラインでは
top-$k$に含まれるが、提案手法のマルチスケールフィルタリングにより
自動的に下位にランクされる。

具体的には、$\{B, G\}$は$W{=}10$のスケールでのみ密集区間$[45, 72]$を
形成するが、$W{=}20$や$W{=}40$では密度が閾値を下回る。
これは粒度ゆらぎによる局所的な偶然の集中であり、
真の時間的構造を反映していないことをスケール空間分析が正しく識別している。
対照的に、真のパターン$\{A, B\}$は$W{=}10$で$[120, 178]$、
$W{=}20$で$[115, 182]$、$W{=}40$で$[110, 190]$と
スケールを超えて一貫した密集区間を維持しており、
リッジ幅4がこの安定性を数値的に捉えている。

% ==============================================================
\section{考察}\label{sec:discussion}
% ==============================================================

\subsection{構造認識型マイニングの意義}

本フレームワークは、従来の密集区間マイニングにおける2つの構造的盲点を
同時に解消する。第1に、二分スケール階層により、
$[W_0, 2^L W_0]$内の任意のスケールで密集であるパターンが、
ユーザが正しいウィンドウサイズを推測する必要なく捕捉される。
第2に、スケール空間密集リッジにより、トランザクション粒度のゆらぎや
セッション境界の不整合に起因するスプリアスパターンが自動的にフィルタリングされる。
E5の結果が示すように、Pre-DedupやSession-Splitなどの前処理ベースラインは
アドホックな構造仮定に依存するが、提案手法はスケール不変性という原理的な
基準により構造認識を実現する。

\subsection{実用上の考慮事項}

$W_0$はパラメータとして残るが、実験結果はその選択に対して結果が
非感受的であることを示している。実用上、$W_0$は小さな値（例えば5--10）に
設定でき、二分階層が関連する範囲を自動的にカバーする。
閾値$\theta_0$はtop-$k$定式化により実効的に排除され、
結果セットに入るために必要な最小スコアがアルゴリズム内部で決定される。

\subsection{制約事項}

現在の実装は数千トランザクションまでのデータセットに適したPythonプロトタイプを
使用している。より大規模な応用には、本フレームワークの2段階開発プロセスに
従ったRust実装が必要となる。
二分スケール階層は2のべき乗のウィンドウサイズのみをカバーする；
非二分の拡張は計算コストの増加と引き換えに、より細かいスケール解像度を
提供しうる。

% ==============================================================
\section{結論}\label{sec:conclusion}
% ==============================================================

本論文では、従来の$(W, \theta)$パラメータを直感的な単一パラメータ$k$に置き換え、
さらにデータの時間的構造を原理的に考慮する、
密集区間マイニングのための構造認識型フレームワークを提示した。
本手法はパターンのランキングのためのDense Coverage Score、
マルチスケール分析のための二分スケール階層、
およびスケール不変なパターン特定のためのスケール空間密集リッジを導入した。
反単調性枝刈りを持つ分枝限定法アルゴリズムにより、
効率的な計算で正確なtop-$k$結果を保証する。

実験では、網羅的グリッドサーチとの品質同等性、前処理ベースライン
（Pre-Dedup, Session-Split）に対するスプリアスパターン削減、
効果的な枝刈り、ほぼ線形のスケーラビリティ、
およびロバストなパラメータ非感受性を実証した。
パターンレベルのケーススタディ（E6）により、
スケール空間密集リッジが真のパターンとスプリアスパターンを
スケール持続性の観点から原理的に区別できることを示した。
今後の研究では、インクリメンタルMSDCS更新によるストリーミング設定への拡張、
および小売トランザクションログやIoTイベントストリームなどの大規模実データセットへの
適用を予定している。

% ==============================================================
\bibliographystyle{IEEEtran}
\bibliography{refs}

\end{CJK}
\end{document}
